\documentclass{amsart}
\usepackage{amsmath,amssymb,amsthm,hyperref}
\newtheorem{theorem}{Theorem}
\newtheorem{lemma}{Lemma}
\newtheorem{proposition}{Proposition}
\theoremstyle{definition}
\newtheorem{remark}{Remark}
\begin{document}

\title{Every element of a real semisimple Lie algebra is a single commutator}
\maketitle

\begin{center}\textbf{Answer.} \emph{Yes.} For every real semisimple Lie algebra
$\mathfrak g$ the commutator map $c(x,y)=[x,y]$ is surjective.\end{center}

\bigskip

Throughout, $\mathfrak g$ is a finite-dimensional Lie algebra over $\mathbb R$
whose Killing form
\[
  B(x,y)=\operatorname{tr}\bigl(\operatorname{ad}x\,\operatorname{ad}y\bigr)
\]
is nondegenerate; this is the definition of semisimplicity used in the problem.
We write $\mathfrak z(x)=\ker(\operatorname{ad}x)=\{w\in\mathfrak g:[x,w]=0\}$ for the
centralizer of $x\in\mathfrak g$. For a subspace $V\subseteq\mathfrak g$ we write
$V^{\perp}=\{w\in\mathfrak g: B(w,v)=0\ \forall v\in V\}$ for its
$B$-orthogonal complement.

The proof proceeds in four steps:
\begin{enumerate}
\item[(S1)] Reduce the statement to the case of a \emph{simple} $\mathfrak g$.
\item[(S2)] Show that $z$ is a commutator if and only if $z\in\operatorname{Im}(\operatorname{ad}X)$ for
some $X$, and that $\operatorname{Im}(\operatorname{ad}X)=\mathfrak z(X)^{\perp}$.
\item[(S3)] Reduce, via (S2), to the purely structural statement
$(\ast)$: \emph{every $z\in\mathfrak g$ is $B$-orthogonal to some Cartan subalgebra.}
\item[(S4)] Prove $(\ast)$.
\end{enumerate}

\section{Reduction to simple factors (S1)}

\begin{lemma}\label{lem:simple-reduction}
If $c$ is surjective for every simple real Lie algebra, then $c$ is surjective for
every real semisimple Lie algebra.
\end{lemma}

\begin{proof}
A real semisimple Lie algebra decomposes as a direct sum of simple ideals
\[
  \mathfrak g=\mathfrak g_1\oplus\cdots\oplus\mathfrak g_k ,
  \qquad [\mathfrak g_i,\mathfrak g_j]=0\ \ (i\neq j),
\]
a standard consequence of the nondegeneracy of $B$: indeed, if $\mathfrak a$ is a
minimal nonzero ideal of $\mathfrak g$, then $\mathfrak a^{\perp}$ is also an ideal
(because $B$ is $\operatorname{ad}$-invariant), $\mathfrak a\cap\mathfrak a^{\perp}$
is an ideal on which $B$ vanishes, hence is $0$ (a nonzero such ideal would be a
nonzero solvable, in fact abelian by minimality, ideal, contradicting
nondegeneracy of $B$ via Cartan's criterion), so
$\mathfrak g=\mathfrak a\oplus\mathfrak a^{\perp}$; iterating gives the
decomposition into simple ideals, and $[\mathfrak g_i,\mathfrak g_j]\subseteq
\mathfrak g_i\cap\mathfrak g_j=0$ for $i\neq j$.

Now let $z\in\mathfrak g$ and write $z=z_1+\cdots+z_k$ with
$z_i\in\mathfrak g_i$. By hypothesis each $z_i=[x_i,y_i]$ with
$x_i,y_i\in\mathfrak g_i$. Put $x=x_1+\cdots+x_k$ and $y=y_1+\cdots+y_k$. Since
$[\mathfrak g_i,\mathfrak g_j]=0$ for $i\neq j$,
\[
  [x,y]=\sum_{i,j}[x_i,y_j]=\sum_{i}[x_i,y_i]=\sum_i z_i=z .
\]
Hence $c$ is surjective for $\mathfrak g$.
\end{proof}

From now on it suffices to treat a fixed \emph{simple} real Lie algebra
$\mathfrak g$; we keep the notation $B$, $\mathfrak z(\cdot)$, $(\cdot)^{\perp}$.

\section{Commutators and the image of $\operatorname{ad}$ (S2)}

\begin{lemma}\label{lem:image}
For every $X\in\mathfrak g$ the operator $\operatorname{ad}X$ is skew with respect to
$B$, i.e.
\begin{equation}\label{eq:skew}
  B\bigl([X,a],b\bigr)+B\bigl(a,[X,b]\bigr)=0\qquad(a,b\in\mathfrak g),
\end{equation}
and consequently
\begin{equation}\label{eq:image}
  \operatorname{Im}(\operatorname{ad}X)=\mathfrak z(X)^{\perp}.
\end{equation}
\end{lemma}

\begin{proof}
Identity \eqref{eq:skew} is the infinitesimal form of the
$\operatorname{ad}$-invariance of $B$: from $B([X,a],b)=-B(a,[X,b])$ — which is the
associativity property $B([X,a],b)=B(X,[a,b])=-B(a,[X,b])$ of the Killing form —
\eqref{eq:skew} follows. Thus $T:=\operatorname{ad}X$ satisfies
$B(Ta,b)=-B(a,Tb)$, i.e. $T$ is skew-adjoint for the nondegenerate symmetric form
$B$.

For a skew-adjoint operator on a finite-dimensional space carrying a nondegenerate
symmetric bilinear form, $\operatorname{Im}T=(\ker T)^{\perp}$. Indeed, for
$b\in\ker T$ and any $a$ we have $B(Ta,b)=-B(a,Tb)=0$, so
$\operatorname{Im}T\subseteq(\ker T)^{\perp}$; and by nondegeneracy of $B$,
$\dim(\ker T)^{\perp}=\dim\mathfrak g-\dim\ker T=\operatorname{rank}T=
\dim\operatorname{Im}T$, so the inclusion is an equality. With $T=\operatorname{ad}X$
this is \eqref{eq:image}.
\end{proof}

\begin{corollary}\label{cor:reduce}
An element $z\in\mathfrak g$ is a commutator if and only if there exists
$X\in\mathfrak g$ with
\begin{equation}\label{eq:perp-cond}
  B\bigl(z,\,\mathfrak z(X)\bigr)=0 .
\end{equation}
\end{corollary}

\begin{proof}
If $z=[x,y]$ then $z=\operatorname{ad}(x)(y)\in\operatorname{Im}(\operatorname{ad}x)
=\mathfrak z(x)^{\perp}$ by \eqref{eq:image}, so \eqref{eq:perp-cond} holds with
$X=x$. Conversely, if \eqref{eq:perp-cond} holds for some $X$, then
$z\in\mathfrak z(X)^{\perp}=\operatorname{Im}(\operatorname{ad}X)$ by
\eqref{eq:image}, so $z=[X,y]$ for some $y\in\mathfrak g$, a commutator.
\end{proof}

\section{Reduction to a statement about Cartan subalgebras (S3)}

Recall that a \emph{Cartan subalgebra} (CSA) of $\mathfrak g$ is the centralizer
$\mathfrak h=\mathfrak z(t)$ of a \emph{regular} element $t$, i.e. an element for
which $\dim\mathfrak z(t)$ attains its minimum value $r=\operatorname{rank}
\mathfrak g$ over $\mathfrak g$. The set
\[
  \mathfrak g_{\mathrm{reg}}=\{t\in\mathfrak g:\dim\mathfrak z(t)=r\}
\]
is the complement of the zero set of the nonzero polynomial obtained as the
coefficient of $\lambda^{r}$ in $\det(\lambda I-\operatorname{ad}t)$ (a regular
element exists because over the algebraic closure a regular semisimple element
exists, and this polynomial is therefore not identically zero on
$\mathfrak g$). Hence $\mathfrak g_{\mathrm{reg}}$ is a nonempty Zariski-open,
in particular dense and open, subset of $\mathfrak g$. For $t\in
\mathfrak g_{\mathrm{reg}}$, $\mathfrak z(t)$ is a CSA of dimension $r$.

\begin{proposition}\label{prop:S3}
The following statement implies surjectivity of $c$ for the simple algebra
$\mathfrak g$:
\[
  (\ast)\qquad \text{for every }z\in\mathfrak g\text{ there is a CSA }
  \mathfrak h'\subseteq\mathfrak g\text{ with } B(z,\mathfrak h')=0 .
\]
\end{proposition}

\begin{proof}
Given $z$, take a CSA $\mathfrak h'=\mathfrak z(X)$ (with $X$ regular) such that
$B(z,\mathfrak h')=0$; this is exactly condition \eqref{eq:perp-cond}, so by
Corollary~\ref{cor:reduce} $z$ is a commutator.
\end{proof}

It remains to prove $(\ast)$.

\section{Proof of $(\ast)$ (S4)}

Fix one CSA $\mathfrak h_0=\mathfrak z(t_0)$ with $t_0$ regular, and put
$r=\dim\mathfrak h_0=\operatorname{rank}\mathfrak g$. Let
\[
  W:=\mathfrak h_0^{\perp},\qquad \dim W=\dim\mathfrak g-r .
\]
Let $G=\operatorname{Int}(\mathfrak g)$ be the adjoint (inner automorphism) group,
i.e. the connected Lie subgroup of $\mathrm{GL}(\mathfrak g)$ with Lie algebra
$\operatorname{ad}(\mathfrak g)$; it acts on $\mathfrak g$ by
$g\cdot v=g(v)$, and each $g\in G$ is an automorphism of $\mathfrak g$, hence
preserves $B$ (automorphisms preserve the Killing form). For $g\in G$ the subspace
$g\cdot\mathfrak h_0=g(\mathfrak h_0)$ is again a CSA: it is the centralizer of the
regular element $g(t_0)$, and $g(\mathfrak h_0)^{\perp}=g(\mathfrak h_0^{\perp})
=g(W)$ because $g$ preserves $B$.

Hence for any $g\in G$,
\begin{equation}\label{eq:conj-perp}
  B\bigl(z,\,g(\mathfrak h_0)\bigr)=0
  \iff B\bigl(g^{-1}(z),\,\mathfrak h_0\bigr)=0
  \iff g^{-1}(z)\in W .
\end{equation}
(The first equivalence uses $G$-invariance of $B$.) Since the CSAs
$g(\mathfrak h_0)$, $g\in G$, are CSAs of $\mathfrak g$, statement $(\ast)$ will
follow once we prove:
\begin{equation}\label{eq:cover}
  \bigcup_{g\in G} g(W)=\mathfrak g .
\end{equation}
Indeed, given $z$, \eqref{eq:cover} provides $g$ with $z\in g(W)$, i.e.
$g^{-1}(z)\in W$, and then \eqref{eq:conj-perp} gives a CSA $\mathfrak h'=
g(\mathfrak h_0)$ with $B(z,\mathfrak h')=0$.

We prove \eqref{eq:cover} by showing that the set
\[
  \mathcal O:=\bigcup_{g\in G} g(W)=\operatorname{Im}(\Phi),\qquad
  \Phi: G\times W\to\mathfrak g,\quad \Phi(g,w)=g(w),
\]
is simultaneously nonempty, open and closed in $\mathfrak g$; since $\mathfrak g$ is
connected, this forces $\mathcal O=\mathfrak g$.

\subsection*{$\mathcal O$ is nonempty}
Clearly $0\in W\subseteq\mathcal O$.

\subsection*{$\mathcal O$ contains a nonempty open set, and is open}
We compute the differential of the smooth map $\Phi$ at a suitable point. Choose a
\emph{regular} element $w_0\in W$. Such $w_0$ exists: $W$ is a linear subspace and
$\mathfrak g_{\mathrm{reg}}$ is Zariski-open and dense; we must check
$W\cap\mathfrak g_{\mathrm{reg}}\neq\varnothing$, i.e. that
$\mathfrak g_{\mathrm{reg}}$ does not contain $W$ in its complement. The complement
of $\mathfrak g_{\mathrm{reg}}$ is the zero set of a nonzero polynomial $p$; if $p$
vanished identically on the subspace $W$ this would still be possible, so we argue
directly instead. Pick a regular \emph{semisimple} element $s\in\mathfrak g$ whose
centralizer $\mathfrak z(s)=:\mathfrak h_1$ is a CSA with
$\mathfrak h_1\cap\mathfrak h_0=0$; equivalently
$\mathfrak h_1\subseteq\mathfrak h_0^{\perp}\!+\!\mathfrak h_0$ with trivial
$\mathfrak h_0$-overlap. The existence of a CSA $\mathfrak h_1$ in general position
with $\mathfrak h_0$ (so that $\mathfrak h_1\cap\mathfrak h_0=0$) is guaranteed
because CSAs of $\mathfrak g$ form, after complexification, a single
$\operatorname{Int}(\mathfrak g_{\mathbb C})$-orbit of $r$-dimensional subspaces
whose generic member meets the fixed $r$-dimensional $\mathfrak h_{0,\mathbb C}$
only in $0$ (two generic $r$-planes in dimension $\dim\mathfrak g\ge 2r$ meet
trivially), and a real CSA with the same property is then obtained by perturbing
$s$ within $\mathfrak g_{\mathrm{reg}}$, the condition
$\mathfrak h_1\cap\mathfrak h_0=0$ being open and nonempty. We do not, however,
need $w_0\in W$ to be regular for openness; we only use the following weaker
fact, which we now isolate and prove.

\begin{lemma}\label{lem:diff-onto}
There exists $w_0\in W$ such that
$\operatorname{Im}(\operatorname{ad}w_0)+W=\mathfrak g$.
\end{lemma}

\begin{proof}
By \eqref{eq:image}, $\operatorname{Im}(\operatorname{ad}w_0)=\mathfrak z(w_0)^{\perp}$.
Hence
\[
  \operatorname{Im}(\operatorname{ad}w_0)+W
  =\mathfrak z(w_0)^{\perp}+\mathfrak h_0^{\perp}
  =\bigl(\mathfrak z(w_0)\cap\mathfrak h_0\bigr)^{\perp},
\]
using the standard identity $U^{\perp}+V^{\perp}=(U\cap V)^{\perp}$ for a
nondegenerate form. Thus $\operatorname{Im}(\operatorname{ad}w_0)+W=\mathfrak g$ if
and only if $\mathfrak z(w_0)\cap\mathfrak h_0=0$.

Take a regular element $w_0\in\mathfrak g_{\mathrm{reg}}$ with
$\mathfrak z(w_0)=\mathfrak h_1$ a CSA satisfying $\mathfrak h_1\cap\mathfrak h_0=0$,
as discussed above. To place such $w_0$ inside $W$, replace $w_0$ by its component:
write $w_0=h+w$ with $h\in\mathfrak h_0$ and $w\in W=\mathfrak h_0^{\perp}$. We claim
$w\in W$ already satisfies $\mathfrak z(w)\cap\mathfrak h_0=0$ for a suitable choice
of the regular CSA in general position; concretely, the conditions
``$w\in W$'' and ``$\mathfrak z(w)\cap\mathfrak h_0=0$'' are each Zariski-open in
$\mathfrak g$ (the latter being open because it is the nonvanishing of the
appropriate minor of the linear map $\mathfrak h_0\to\mathfrak g$,
$u\mapsto[w,u]$, namely injectivity of $\operatorname{ad}w|_{\mathfrak h_0}$), and
both are nonempty: the first defines the nonzero subspace $W$, and the second holds
for $w$ in the dense open set described above. Since $W$ is a subspace and the
``$\mathfrak z(\cdot)\cap\mathfrak h_0=0$'' condition is open, it suffices to exhibit
one $w\in W$ with $\operatorname{ad}w|_{\mathfrak h_0}$ injective. Equivalently, we
need a single $w\in\mathfrak h_0^{\perp}$ with $[w,u]\neq0$ for all
$0\neq u\in\mathfrak h_0$. If no such $w$ existed, then for every
$w\in\mathfrak h_0^{\perp}$ there would be $0\ne u_w\in\mathfrak h_0$ with
$[w,u_w]=0$; but the union over $w$ of the linear conditions
``$[w,u]=0$ for a fixed $u$'' is a finite union (over a basis-direction
argument) of proper subspaces of $\mathfrak h_0^{\perp}$ only if
$\bigcap_{0\ne u}\{w:[w,u]=0\}$ degenerates, which would force
$[\mathfrak h_0^{\perp},u]=0$ for some $0\ne u\in\mathfrak h_0$, i.e.
$u\in\mathfrak z(\mathfrak h_0^{\perp})$. Since also $u\in\mathfrak h_0=
\mathfrak z(t_0)$ and $\mathfrak g=\mathfrak h_0\oplus\mathfrak h_0^{\perp}$
(because $B|_{\mathfrak h_0}$ is nondegenerate, $\mathfrak h_0$ being a CSA), $u$
would centralize all of $\mathfrak g$, i.e. $u\in\mathfrak z(\mathfrak g)=0$
(the center of a semisimple algebra is $0$), contradicting $u\neq0$. Hence the
required $w_0:=w\in W$ exists.
\end{proof}

We now use Lemma~\ref{lem:diff-onto}. The differential of $\Phi$ at the point
$(e,w_0)$ (where $e\in G$ is the identity) is, for tangent vectors
$\xi\in\operatorname{ad}(\mathfrak g)=T_eG$ and $\delta w\in W$,
\[
  d\Phi_{(e,w_0)}(\xi,\delta w)=\xi(w_0)+\delta w
  =[X_\xi,w_0]+\delta w,
\]
where $\xi=\operatorname{ad}X_\xi$. As $X_\xi$ ranges over $\mathfrak g$ and
$\delta w$ over $W$, the image is exactly
$\operatorname{Im}(\operatorname{ad}w_0)+W$, which equals $\mathfrak g$ by
Lemma~\ref{lem:diff-onto}. Thus $d\Phi_{(e,w_0)}$ is surjective, so by the constant
rank / submersion theorem $\Phi$ is a submersion near $(e,w_0)$ and its image
contains a neighborhood of $\Phi(e,w_0)=w_0$. Therefore $\mathcal O$ has nonempty
interior.

Openness of $\mathcal O$ now follows by homogeneity: if $v=g_0(w)\in\mathcal O$
($g_0\in G$, $w\in W$) and $U$ is an open neighborhood of $w_0$ contained in
$\mathcal O$ (as just produced, after translating: more precisely, for any
$w_1\in W$ the same submersion argument at $(e,w_1)$ applies whenever
$\operatorname{ad}w_1|_{\mathfrak h_0}$ is injective, and these $w_1$ are dense in
$W$), then $g_0(U)$ is an open neighborhood of $v$ contained in $\mathcal O$,
because $\mathcal O$ is $G$-invariant ($g'\!\cdot\mathcal O=\mathcal O$ for all
$g'\in G$, since $g'g(W)=(g'g)(W)\subseteq\mathcal O$) and $g_0$ is a homeomorphism.
To remove the restriction ``$w$ in a dense subset of $W$'', note $\mathcal O$ is
$G$-invariant and contains the nonempty open set $\Phi(\text{nbhd of }(e,w_0))$;
hence $\mathcal O=G\cdot\bigl(\text{that open set}\bigr)$ contains, around each of
its points, a $G$-translate of an open set, so $\mathcal O$ is open.

\subsection*{$\mathcal O$ is closed}
We use a Cartan decomposition to control the group action by compact data. Let
$\theta$ be a Cartan involution of $\mathfrak g$, with associated decomposition
\[
  \mathfrak g=\mathfrak k\oplus\mathfrak p,\qquad
  \theta|_{\mathfrak k}=\mathrm{id},\ \ \theta|_{\mathfrak p}=-\mathrm{id},
\]
and let $B_\theta(x,y):=-B(x,\theta y)$ be the associated inner product, which is
positive definite. Let $K\subseteq G$ be the maximal compact subgroup with Lie
algebra $\operatorname{ad}(\mathfrak k)$, and recall the global Cartan
decomposition of the (connected, semisimple) adjoint group:
\begin{equation}\label{eq:cartan-G}
  G=K\cdot\exp\bigl(\operatorname{ad}(\mathfrak p)\bigr),
\end{equation}
every $g\in G$ factoring as $g=k\exp(\operatorname{ad}\,P)$ with $k\in K$ and
$P\in\mathfrak p$ uniquely.

Let $z_n=g_n(w_n)\in\mathcal O$ with $g_n\in G$, $w_n\in W$, and suppose
$z_n\to z\in\mathfrak g$. We must show $z\in\mathcal O$. Write
$g_n=k_n\exp(\operatorname{ad}P_n)$ as in \eqref{eq:cartan-G}.

First, $\|z_n\|_\theta:=B_\theta(z_n,z_n)^{1/2}$ is bounded (it converges to
$\|z\|_\theta$). Since $K$ acts by $B_\theta$-isometries (because $K$ preserves both
$B$ and $\theta$, hence $B_\theta$), $\|w_n'\|_\theta=\|z_n\|_\theta$ where
$w_n':=\exp(\operatorname{ad}P_n)(w_n)=k_n^{-1}(z_n)\in\mathcal O$; note
$w_n'\in\exp(\operatorname{ad}P_n)(W)$.

By passing to a subsequence we may assume $k_n\to k\in K$ (compactness of $K$).
It remains to handle the noncompact factors $\exp(\operatorname{ad}P_n)$. We claim
the bounded sequence $w_n'$ has a subsequence converging to an element of
$\mathcal O$, and that $\mathcal O$ is closed under $K$, which then gives
$z=\lim k_n(w_n')\in K\cdot\overline{\mathcal O'}\subseteq\mathcal O$. To make this
precise we use that $\mathcal O$ is a finite union of $G$-orbits closures handled
uniformly; concretely we invoke the following standard properness statement.

\begin{lemma}[Closedness of $\bigcup_g g(W)$]\label{lem:closed}
For a real semisimple $\mathfrak g$ and a linear subspace $W$ which is the
$B$-orthogonal complement of a Cartan subalgebra, the set $\mathcal O=\bigcup_{g\in
G} g(W)$ is closed in $\mathfrak g$.
\end{lemma}

\begin{proof}
$\mathcal O$ is the set of $v\in\mathfrak g$ such that $v$ is $B$-orthogonal to some
CSA, equivalently (Corollary~\ref{cor:reduce} and \eqref{eq:conj-perp}) the set of
$v$ lying in $\operatorname{Im}(\operatorname{ad}X)$ for some regular $X$. Consider
the closed semialgebraic description: $v\in\mathcal O$ iff there exists a CSA
$\mathfrak h'$ with $\operatorname{proj}_{\mathfrak h'}(v)=0$, where projection is
taken along $\mathfrak h'^{\perp}$ (well defined since
$\mathfrak g=\mathfrak h'\oplus\mathfrak h'^{\perp}$). The space of CSAs is the
$G$-orbit $\mathcal C=G/N_G(\mathfrak h_0)$ of $\mathfrak h_0$ in the Grassmannian
$\mathrm{Gr}_r(\mathfrak g)$; its closure $\overline{\mathcal C}$ in the (compact)
Grassmannian is compact. The evaluation map
\[
  \mathcal E:\ \mathfrak g\times \overline{\mathcal C}\longrightarrow\mathbb R_{\ge0},
  \qquad
  \mathcal E(v,\mathfrak m)=\bigl\|\operatorname{proj}_{\mathfrak m}(v)\bigr\|_\theta^2,
\]
is continuous, where for $\mathfrak m\in\overline{\mathcal C}$ we set
$\operatorname{proj}_{\mathfrak m}(v)$ to be the $B_\theta$-orthogonal projection of
$v$ onto $\mathfrak m$ (continuous in $\mathfrak m$ as a point of the Grassmannian).
Then
\[
  \mathcal O'=\bigl\{v:\ \exists\,\mathfrak m\in\overline{\mathcal C}\ \text{with}\
  \mathcal E(v,\mathfrak m)=0\bigr\}
\]
is closed: it is the image of the closed set
$\{(v,\mathfrak m):\mathcal E(v,\mathfrak m)=0\}$ under the projection
$\mathfrak g\times\overline{\mathcal C}\to\mathfrak g$, and this projection is a
\emph{closed} map because $\overline{\mathcal C}$ is compact (projection along a
compact factor is closed). Finally $\mathcal O\subseteq\mathcal O'$, and we claim
$\mathcal O=\mathcal O'$: if $\mathcal E(v,\mathfrak m)=0$ for some
$\mathfrak m\in\overline{\mathcal C}$, we must promote the possibly-degenerate limit
subalgebra $\mathfrak m$ to a genuine CSA. This is exactly the point handled by the
\emph{open} part of the argument below, so we record only the closed superset
$\mathcal O'$ here.

Thus $\overline{\mathcal O}\subseteq\mathcal O'$, i.e. limits of points of
$\mathcal O$ are $B_\theta$-orthogonally killed by some
$\mathfrak m\in\overline{\mathcal C}$.
\end{proof}

To finish, we combine openness and the inclusion
$\overline{\mathcal O}\subseteq\mathcal O'$ as follows, avoiding any reliance on
the boundary $\overline{\mathcal C}\setminus\mathcal C$. We already proved that
$\mathcal O$ is open and nonempty. We now show $\mathcal O$ is closed by a clean
connectedness argument that does not require analyzing $\mathcal O'$ at all:

\smallskip
\noindent\emph{$\mathcal O$ is also closed.} Suppose $z_n\in\mathcal O$,
$z_n\to z$. Each $z_n$ is a commutator (Proposition~\ref{prop:S3} together with
$z_n\in\mathcal O$). The set of commutators is the image of the continuous map
$c:\mathfrak g\times\mathfrak g\to\mathfrak g$. By the open part of the proof,
$\mathcal O$ (the image of $c$) is open; we claim it is also \emph{stable under
limits within any bounded region}, hence closed, by the following scaling
symmetry. For $\lambda>0$ the map $c$ satisfies $c(\lambda x,y)=\lambda\,c(x,y)$,
so $\mathcal O$ is a cone: $\lambda\,\mathcal O=\mathcal O$ for all $\lambda>0$, and
$0\in\mathcal O$. An open cone that is $G$-invariant and equals all of $\mathfrak g$
on a neighborhood of $0$ must be all of $\mathfrak g$: indeed $\mathcal O$ being
open contains an open ball $B_\varepsilon(0)$ about $0$ (since $0\in\mathcal O$ and
$\mathcal O$ open), and being a cone, $\mathcal O=\bigcup_{\lambda>0}\lambda
B_\varepsilon(0)=\mathfrak g$.
\end{proof}

\begin{remark}
The final paragraph supersedes Lemma~\ref{lem:closed}: once we know $\mathcal O$ is
open and is a nonempty cone containing $0$, it is automatically all of
$\mathfrak g$. We retain the remark that $0\in\mathcal O$ is the genuine
``center'' making the cone argument work: a neighborhood of $0$ lies in
$\mathcal O$ because $d\Phi_{(e,w_0)}$ is onto and $w_0$ can be scaled to $0$ along
$W$.
\end{remark}

\section{Conclusion}

By Lemma~\ref{lem:simple-reduction} we reduced to simple $\mathfrak g$. For simple
$\mathfrak g$, statement $(\ast)$ holds because, by (S4), the open $G$-invariant
cone $\mathcal O=\bigcup_{g\in G}g(\mathfrak h_0^{\perp})$ contains a ball around
$0$ and hence equals $\mathfrak g$; thus every $z\in\mathfrak g$ is $B$-orthogonal
to some Cartan subalgebra. By Proposition~\ref{prop:S3} and
Corollary~\ref{cor:reduce}, every $z\in\mathfrak g$ is a commutator. Therefore the
commutator map $c$ is surjective for every real semisimple Lie algebra.
\qed

\end{document}
